% ==============================================================================
% SECCIÓN 6: ARQUITECTURA FUNCIONAL, MÓDULOS Y MODELO DE DATOS
% ==============================================================================

\section{Arquitectura Funcional y Alcance de Módulos del SIGC-CUSCO}

\subsection{Estructura Modular del Sistema}
El sistema se organiza en ocho módulos funcionales coordinados (Tabla~\ref{tab:modulos}).

\begin{table}[H]
\centering
\small
\renewcommand{\arraystretch}{1.18}
\begin{tabularx}{\textwidth}{P{3.2cm}P{5.0cm}Y}
\toprule
\textbf{Módulo} & \textbf{Responsabilidad Principal} & \textbf{Operaciones Clave} \\
\midrule
\textbf{1. Usuarios y Roles} & Autenticación y control de acceso (RBAC). & Login seguro, roles: Admin, Docente, Alumno, Institución. \\
\midrule
\textbf{2. Instituciones} & Gestión multitenant de entidades organizadoras. & Registro de municipalidades, empresas de turismo y sellos. \\
\midrule
\textbf{3. Capacitaciones} & Catálogo y parametrización de cursos. & Creación de cursos, fechas, horas, cupos y docentes. \\
\midrule
\textbf{4. Inscripciones} & Registro de participantes y cupos. & Formulario público, validación de DNI y confirmación por correo. \\
\midrule
\textbf{5. Asistencia} & Control digital y auditable por sesión. & Marcado mediante código QR / PIN con soporte \textit{offline}. \\
\midrule
\textbf{6. Evaluaciones} & Registro de notas y actas académicas. & Calificaciones (0--20) y cálculo automático de aprobación. \\
\midrule
\textbf{7. Certificados} & Emisión y validación documental. & Generación masiva de PDFs con firma digital, QR y SHA-256. \\
\midrule
\textbf{8. Reportes} & Analítica institucional para auditoría. & Asistencia, índice de aprobación y exportación oficial a Excel. \\
\bottomrule
\end{tabularx}
\caption{Módulos funcionales del sistema SIGC-CUSCO.}
\label{tab:modulos}
\end{table}

\subsection{Flujo Operativo del Proceso Formativo}
El ciclo de vida de una capacitación en el sistema sigue una secuencia verificable y sin fricciones:

\begin{center}
\begin{tikzpicture}[
    stepnode/.style={rectangle, draw=blue!70!black, fill=blue!5, text width=12cm, rounded corners=3pt, inner sep=4pt, font=\small},
    arr/.style={draw=blue!70!black, -Latex, thick}
]
\node[stepnode] (s1) {\textbf{1. Planificación:} La institución crea el curso, fija vacantes, fechas y asigna al capacitador.};
\node[stepnode, below=0.20cm of s1] (s2) {\textbf{2. Inscripción Digital:} Los participantes se registran en línea con validación de DNI en tiempo real.};
\node[stepnode, below=0.20cm of s2] (s3) {\textbf{3. Control de Asistencia:} Marcado ágil mediante código QR en cada sesión (soporte offline-first).};
\node[stepnode, below=0.20cm of s3] (s4) {\textbf{4. Evaluación Académica:} El docente califica exámenes y el sistema determina automáticamente la aprobación.};
\node[stepnode, below=0.20cm of s4] (s5) {\textbf{5. Certificación Automática:} Generación masiva de diplomas PDF con código QR y hash SHA-256.};
\node[stepnode, below=0.20cm of s5] (s6) {\textbf{6. Verificación Pública y Analítica:} Consulta web del certificado y generación de reportes para auditoría.};

\draw[arr] (s1) -- (s2);
\draw[arr] (s2) -- (s3);
\draw[arr] (s3) -- (s4);
\draw[arr] (s4) -- (s5);
\draw[arr] (s5) -- (s6);
\end{tikzpicture}
\end{center}

\subsection{Modelo de Datos Esencial (Entidades Principales)}
\begin{itemize}\setlength\itemsep{0.15em}\raggedright
    \item \textbf{Institución:} \texttt{id\_institucion}, \texttt{ruc}, \texttt{razon\_social}, \texttt{tipo}, \texttt{logo\_url}, \texttt{estado}.
    \item \textbf{Usuario:} \texttt{id\_usuario}, \texttt{dni}, \texttt{nombres}, \texttt{apellidos}, \texttt{correo}, \texttt{password\_hash}, \texttt{rol}.
    \item \textbf{Capacitación:} \texttt{id\_capacitacion}, \texttt{id\_institucion}, \texttt{id\_docente}, \texttt{titulo}, \texttt{horas}, \texttt{cupos}.
    \item \textbf{Inscripción:} \texttt{id\_inscripcion}, \texttt{id\_capacitacion}, \texttt{id\_participante}, \texttt{fecha\_registro}, \texttt{estado}.
    \item \textbf{Asistencia:} \texttt{id\_asistencia}, \texttt{id\_inscripcion}, \texttt{fecha\_sesion}, \texttt{estado}, \texttt{metodo}.
    \item \textbf{Evaluación:} \texttt{id\_evaluacion}, \texttt{id\_inscripcion}, \texttt{nota\_final} (0 a 20), \texttt{condicion}.
    \item \textbf{Certificado:} \texttt{id\_certificado}, \texttt{codigo\_unico}, \texttt{hash\_sha256}, \texttt{url\_pdf}, \texttt{url\_qr}.
\end{itemize}

\subsection{Plan de Escalabilidad en 5 Fases}
\begin{enumerate}\setlength\itemsep{0.1em}
    \item \textbf{Fase 1 (MVP Local):} Municipalidades de Cusco (capacitaciones en SIGA, OSCE y residuos sólidos).
    \item \textbf{Fase 2 (Sector Turismo):} Hoteles, restaurantes y agencias de viaje (atención al cliente y calidad turística).
    \item \textbf{Fase 3 (Educación Continua):} Institutos técnicos y colegios profesionales del Cusco.
    \item \textbf{Fase 4 (Emprendedores y Gremios):} Asociaciones comerciales, artesanales y PYMEs locales.
    \item \textbf{Fase 5 (Analítica con IA):} Detección de brechas formativas laborales y validación descentralizada.
\end{enumerate}
